\begin{abstract}
Establishing a permanent human presence on the Moon requires closing habitat life-support loops through Biogenerative Life Support Systems (BLSS). Terrestrial resupply imposes prohibitive launch costs, while direct agricultural cultivation in unrefined lunar regolith is severely constrained by submicroscopic nanophase metallic iron ($\mathrm{npFe}^0$), abiotic generation of reactive oxygen species (ROS) via modified Fenton cascades, broken silicate dangling bonds, and sharp mineral shards that induce physical shearing and plant root apical growth arrest. Furthermore, fractional lunar gravity ($1/6\,g$) alters the Bond number, resulting in capillary entrapment, erratic nutrient transport, and pervasive root hypoxia in polydisperse regolith beds. Here, we review an alternative in-situ bio-infrastructure paradigm: transforming unrefined regolith into engineered porous ceramics via thermal sintering. High-temperature processing ($>1050^\circ\mathrm{C}$) passivates reactive geochemical interfaces by solid-solubilizing $\mathrm{npFe}^0$ into stable augite and olivine matrices, rounds grain boundaries, and enables precise control over pore throat diameters ($30\text{--}60\,\mu\mathrm{m}$). This pore architecture stabilizes matric potential ($-2.5$ to $-5.0\,\mathrm{kPa}$), ensuring continuous capillary nutrient wicking while maintaining air-filled porosity above $25\text{--}30\%$. Equivalent System Mass (ESM) lifecycle modeling reveals that the initial launch-mass investment of a sintering plant reaches parity with direct regolith cultivation within $18\text{--}22$ months, driven by $>98.5\%$ water recovery efficiency and the elimination of chemical buffers. When integrated with fungal mycoponics, microbial bioweathering, and cyclic $500^\circ\mathrm{C}$ pyrolytic substrate regeneration, sintered regolith ceramics provide a closed-loop substrate foundation for sustainable extraterrestrial habitation.
\end{abstract}

\vspace{0.5em}
\noindent\textbf{Keywords:} Lunar regolith, Biogenerative Life Support Systems (BLSS), In-Situ Resource Utilization (ISRU), Porous ceramics, Microwave sintering, Nanophase iron, Microgravity hydrodynamics, Equivalent System Mass (ESM).
